\section{Applications}

\DMGtodo[inline]{some possible applications of this result. If the conditions give us intersection numbers, then one should be able to define geodesic currents as a subset of functionals on curves. This allows to give a new definition of Thurston distance in terms of supremums over a subfamily of functionals. Then one can especulate things like: can Teichmueller distance be expressed as supremum over a larger family of functionals on curves?}

\subsection{Redefining Thurston metric}

$d_{\mathit{Th}}(X,Y)=\sup_{f \in \mathcal{C}}\frac{f(X)}{f(Y)}$ (it is not clear if $\mathcal{C}$ here should mean only additive convex functionals (motivation: true for intersection numbers).
\DMGtodo[inline]{What is the smallest family of functionals $\mathcal{C}$ that guarantees this is actually a symmetric distance? }


\subsection{Characterizing binding currents?}

$\operatorname{Syst}(\mu)=0$ if and only if $f(\mu)=0$ for some convex (additive?) functional. (motivation: result by M. Burger on structural theorem for geodesic currents).

\subsection{Anosov representations}

Let 
\[
\partial \Gamma^{[4]}\coloneqq\{(a,b,c,d) \in \partial \Gamma^{4}\colon \{ a,b \} \cap \{ c,d\}=\emptyset \}
\]

\begin{definition}
  A cross-ratio is a continuous function $B\colon \Gamma^{[4]} \to \mathbb{R}$ invariant under the diagonal action of $\pi_1(S)$ and satisfying:
  \begin{itemize}
      \item (symmetry) $B(x,y,z,w)=B(z,w,x,y)$
      \item (additivity) $B(x,y,z,w) + B(x,y,w,u)=B(x,y,z,u)$ for all $x,y,z,w,u \in \partial \Gamma$ such that $(x,y,z,w), (x,y,z,u) \in  \Gamma^{[4]}$
  \end{itemize}
\end{definition}
We say that $B$ is \emph{positive} if for all $x,y,z,w$ in this cyclic order, we have $B(x,y,z,w) \geq 0$. \\
Recall that if $c=[\gamma]$ we can define the length 
\[
  l_{\rho}^{\alpha}(c)=B_{\rho}^{\alpha}(\gamma^+,\gamma^-,z,\gamma(z)).
 \]

The properties of cross ratio ensure this is well-defined.
In \cite{MZ18:PositiveReps} it is shown that if $B_{\rho}^{\alpha}$ is positive, then there exists a unique
geodesic current $\mu_{\rho}^{\alpha}$ so that $l_{\rho}^{\alpha}(c)=i(\mu_{\rho}^{\alpha},c)$. In particular, $l_{\rho}^{\alpha}$ extends to geodesic currents as an intersection number.
We can give another characterization of positivity for cross-ratio in terms of the length function $l_{\rho}^{\alpha}$.
 
\begin{proposition}
 $B_{\rho}^{\alpha}$ is positive if and only if $l_{\rho}^{\alpha}$ satisfies smoothing condition.
\end{proposition}
\begin{proof}
From the above mentioned result in \cite{MZ18:PositiveReps}, the smoothing property follows. \\
On the other hand, if $l_{\rho}^{\alpha}$ satisfies smoothing, since additive union is also clear for $l_{\rho}^{\alpha}$, then we just apply theorem \ref{thm:intersect}
to obtain a geodesic current $\mu_{\rho}^{\alpha}$. Finally, we define the cross-ratio 
\[ 
B_{\rho}^{\alpha}(a,b,c,d)\coloneqq \mu_{\rho}^{\alpha}([a,b)\times[c,d)).
\]

It obviously satisfies the properties of cross-ratio in \cite{MZ18:PositiveReps} because of the definition of geodesic current, and it's positive.

\end{proof}